% Part: lambda-calculus
% Chapter: introduction
% Section: lambda-computable

\documentclass[../../../include/open-logic-section]{subfiles}

\begin{document}

\olfileid{lam}{int}{cmp}
\olsection{\usetoken{S}{lambda definable} অপেক্ষকগুলি গণনসাধ্য}

\begin{thm}
\ollabel{thm:lambda-computable}
কোনো আংশিক অপেক্ষক~$f$-কে একটি ল্যাম্বডা পদ !!{lambda defined} করলে
অপেক্ষকটি গণনসাধ্য।
\end{thm}

\begin{proof}
ধরা যাক, কোনো ল্যাম্বডা পদ~$X$ অপেক্ষক~$f$-কে !!{lambda defined} করে।
$f$ গণনার একটি অনানুষ্ঠানিক পদ্ধতি বর্ণনা করি। নিবেশ $m_0$,
\dots,~$m_{n-1}$ পেলে $X \num{m_0} \ldots \num{m_{n-1}}$ পদটি লেখো।
একটি হ্রাস-বৃক্ষ গড়ো: প্রথমে মূল পদটির সব এক-ধাপের হ্রাস লেখো; তার নিচে
সেগুলির সব এক-ধাপের হ্রাস (অর্থাৎ মূল পদটির দুই-ধাপের হ্রাসগুলি) লেখো;
এবং এভাবেই চালিয়ে যাও। কখনও কোনো সংখ্যাপদে পৌঁছালে সেটিকেই উত্তর হিসেবে
ফেরত দাও; তা না হলে অপেক্ষকটি অসংজ্ঞায়িত।
% BN-SRC-270 TARGET-CORRECTION-BEGIN
\par\smallskip\noindent\textnormal{\small\textbf{উৎস-সংশোধন (BN-SRC-270):} হিমায়িত পাঠের প্রথম অনুসন্ধান-পদে~$\num$ ম্যাক্রোর আর্গুমেন্টদুটি বন্ধনীহীন; ফলে~$m_0$ ও~$m_{n-1}$-এর উপসূচকগুলি সংখ্যাপদ-দাগের বাইরে পড়ে। বাংলা সূত্রে~$\num{m_0}$ ও~$\num{m_{n-1}}$ লেখা হয়েছে, যা নিচের একই পদের সঠিক পুনরাবৃত্তি।} % BN-SRC-270 TARGET-CORRECTION-END

চার্চের থিসিস আহ্বান করলে বলা যায়, এই অপেক্ষকটি গণনসাধ্য। উপপাদ্যটি
প্রমাণের আরও ভালো উপায় হলো এই অনুসন্ধান-পদ্ধতির একটি পুনরাবৃত্ত বর্ণনা
দেওয়া। যেমন, ``$\fn{IsASubterm}$'', ``$\fn{Substitute}$'',
``$\fn{ReducesToInOneStep}$'', ``$\fn{ReductionSequence}$'',
``$\fn{Numeral}$'' ইত্যাদি আদিম পুনরাবৃত্ত অপেক্ষক ও সম্বন্ধের একটি ক্রম
সংজ্ঞায়িত করা যায়। তখন $f(m_0, \dots, m_{n-1})$ গণনার আংশিক পুনরাবৃত্ত
পদ্ধতি হলো $X \num{m_0} \dots \num{m_{n-1}}$ দিয়ে শুরু হয়ে কোনো
সংখ্যাপদে শেষ হয় এমন এক-ধাপের হ্রাসগুলির ক্রম অনুসন্ধান করা, এবং সেই
সংখ্যাপদের অনুরূপ সংখ্যাটি ফেরত দেওয়া। খুঁটিনাটিগুলি দীর্ঘ ও একঘেয়ে,
তবে অন্য দিক দিয়ে নিতান্ত নিয়মমাফিক।
\end{proof}

\end{document}
